% Agent 08: Pages 373--376 (PDF pages 15--17)
% Section 2.6 Radiative Transfer (continued)
% Equations (2.6.11) through (2.6.25)
% Covers: Invariance of I_nu/nu^3, Liouville's theorem,
%         equation of radiative transfer, absorption, stimulated emission,
%         scattering, absorption coefficient kappa_nu, kappa_s, detailed balance

%%% -----------------------------------------------------------------
%%% Invariance of I_nu / nu^3 (continued from earlier on p.~372)
%%% -----------------------------------------------------------------

\textit{Invariance of $I_\nu/\nu^3$}

Choose a particular photon that passes through a given event in spacetime.
Observe that particular photon, and all others in the vicinity, from several different
Lorentz frames at the given event.\ The frequency $\nu$ of the photon will depend
on the choice of Lorentz frame (doppler shift from one frame to another).
The specific intensity $I_\nu$ in the photon's neighborhood will also depend
on the choice of Lorentz frame.\ However, the ratio $I_\nu/\nu^3$ will be Lorentz-invariant.\ (Aside from a factor $h^{-4}$, $I_\nu/\nu^3$ is the invariant number density in
phase space
%
\begin{equation}
\mathscr{N} = \frac{dN}{d^3 x\,d^3 p}
  = \frac{1}{h^3}\,\frac{I_\nu}{\nu^3}\,;
\tag{2.6.11}
\end{equation}
%
see, e.g., \S\,22.6 of MTW.)

When photons propagate freely through curved spacetime (no interaction with
matter), the ratio $I_\nu/\nu^3$ is conserved along the world line of each photon
(Liouville's theorem).

%%% -----------------------------------------------------------------
%%% Equation of Radiative Transfer
%%% -----------------------------------------------------------------

\paragraph*{Equation of Radiative Transfer}

Consider the propagation of photons through a medium (e.g., through gas that
is falling into a black hole).\ Analyze the photon propagation from the viewpoint
of observers at rest in the medium (local rest frame; ``LRF'').\ At each event
denote by $\mathbf{u}$ the 4-velocity of the medium---and hence also of the LRF.\ Focus
attention on all photons in the (phase-space) neighborhood of a given null
geodesic ray.\ Denote by $\mathbf{p}$ the 4-momentum of that ray.\ Then at each event
along the given ray $\mathbf{p}$ is given by
%
\begin{equation}
\mathbf{p} = h\nu(\mathbf{u} + \mathbf{n}).
\tag{2.6.12}
\end{equation}
%
Here $\mathbf{n}$ is a unit vector;
%
\[
\mathbf{n} \cdot \mathbf{u} = 0,
\]
%
and (i) is interpreted in the LRF as the direction of propagation of the ray
(same as 3-vector $\mathbf{n}$ of Figure~2.6.1).\ Also, in eq.~(2.6.12), $\nu$ is the frequency of
any photon which propagates along the ray, and $h\nu$ is its energy, as measured in
the LRF.

Because of the acceleration and shear of the medium, the frequency
$\nu$ of the chosen ray changes from point to point along the ray.\ The change $d\nu$,
when the ray propagates a proper spatial distance $dl$ as seen in the LRF, is
%
\begin{equation}
d\nu = \nu(-\,\mathbf{p} \cdot \mathbf{u}/h)\,dl.
\tag{2.6.13}
\end{equation}
%
(This frame-independent equation is derived easily by geometrical arguments
in the LRF.)\ A straightforward calculation, using the geodesic equation for the
ray
%
\[
\boldsymbol{\nabla}_{\mathbf{p}}\,\mathbf{p} = 0,
\]
%
and using expansion (2.5.17) for $\boldsymbol{\nabla}\mathbf{u}$, reveals
%
\begin{equation}
d\nu/dl = -\nu(\mathbf{n} \cdot \mathbf{a} + \tfrac{1}{3}\theta
  + n^\alpha n^\beta \sigma_{\alpha\beta}).
\tag{2.6.14}
\end{equation}
%
The first term, $-\nu\,\mathbf{n} \cdot \mathbf{a}$, is the ``gravitational redshift'' produced by the acceleration
of the LRF; the second and third terms, $-\nu(\tfrac{1}{3}\theta + n^\alpha n^\beta \sigma_{\alpha\beta})$, are the
``cosmological redshift'' due to the expansion of the medium along the direction
of the ray.

If there were no interaction with the medium, then $I_\nu/\nu^3$ would be conserved
along the ray.\ Hence, the equation of radiative transfer along the given ray must
have the form
%
\begin{equation}
\frac{dI_\nu}{dl} - \left(\frac{3}{\nu}\,\frac{d\nu}{dl}\right) I_\nu
  = \text{(effects of interaction with medium)};
\tag{2.6.15}
\end{equation}
%
i.e.,
%
\[
dI_\nu/dl + (3\mathbf{n} \cdot \mathbf{a} + \theta
  + 3n^\alpha n^\beta \sigma_{\alpha\beta})\,I_\nu
  = \text{(interaction effects)}.
\]
%
Four types of interaction can occur: spontaneous emission of radiation by the
matter; stimulated emission; absorption; and scattering.

We shall assume that on all length scales of interest the medium is isotropic;
and we shall denote its emissivity by
%
\begin{equation}
\varepsilon_\nu \equiv \frac{dE}{d\nu\,dt\,dm_0}
  = \begin{pmatrix} \text{energy emitted spontaneously per unit frequency}\\
    \text{during unit time by a unit rest mass, integrated}\\
    \text{over all angles, and measured in LRF} \end{pmatrix}.
\tag{2.6.16}
\end{equation}
%
(The emissivities for free-free and free-bound transitions were discussed in \S\S\,2.1
and 2.2.)\ Then spontaneous emission contributes
%
\begin{equation}
\left(\frac{dI_\nu}{dl}\right)_{\!\text{spontaneous emission}}
  = \frac{1}{4\pi}\,\rho_0\,\varepsilon_\nu
\tag{2.6.17}
\end{equation}
%
to the specific intensity along the given ray.\ Here $\rho_0$ is the density of rest mass
in the medium.

The rate of absorption and the rate of stimulated emission are both proportional
to the intensity of the passing beam.\ Hence, the effects of absorption and of
stimulated emission can be lumped together into a single \textit{absorption coefficient}
$\kappa_\nu$, defined by
%
\begin{equation}
\left(\frac{dI_\nu}{dl}\right)_{\!\substack{\text{absorption plus}\\[2pt]
  \text{stimulated emission}}}
  = -\rho_0\,\kappa_\nu\,I_\nu.
\tag{2.6.18}
\end{equation}
%
The dimensions of the absorption coefficient are $\mathrm{cm}^2/\mathrm{g}$ (``absorption cross section
per unit mass'' at the given frequency).\ When absorption dominates over
stimulated emission, $\kappa_\nu$ is positive; when stimulated emission dominates,
$\kappa_\nu$ is negative (``negative absorption'').

Scattering is more complicated to treat than emission and absorption.\ Let
%
\begin{equation}
\frac{d\kappa_s}{d\Omega'\,d\nu'}\,(\mathbf{p},\,\mathbf{p}')
\tag{2.6.19}
\end{equation}
%
be the differential cross section (as measured in the LRF) for a unit rest mass to
convert a photon of momentum $\mathbf{p}$ into a photon of momentum $\mathbf{p}'$; and let
%
\[
\kappa_s(\nu) = \int \frac{d\kappa_s}{d\Omega'\,d\nu'}\,d\Omega'\,d\nu'
\]
%
be the total scattering cross-section (``scattering opacity'') per unit rest mass.
For example, in the case of electron scattering by a nonrelativistic plasma
($h\nu \ll m_e c^2$; $kT \ll m_e c^2$; see \S\,4.4), when one neglects the tiny change in photon
frequency the cross sections per unit rest mass are
%
\begin{equation}
\frac{d\kappa_s}{d\Omega'\,d\nu'}\,(\mathbf{p},\,\mathbf{p}')
  = \frac{f_e}{2m_p}\,\sigma_T\,\tfrac{3}{8\pi}\,
    [1 + (\mathbf{n} \cdot \mathbf{n}')^2]\,\delta(\nu' - \nu)
\tag{2.6.20}
\end{equation}
%
\[
\kappa_s = (f_e/m_p)\,\sigma_T = (0.40\;\mathrm{cm}^2/\mathrm{g})\,f_e
\]
%
(Recall: $n = \rho/h\nu$, $n' = \rho'/h\nu'$; $f_e$ is the number of free electrons per baryon in
the plasma.)\ Scattering, as described by the appropriate differential cross section,
can increase $I_\nu$ (scattering into the beam from other directions), or can decrease
$I_\nu$ (scattering out of the beam into other directions):
%
\begin{equation}
\left(\frac{dI_\nu}{dl}\right)_{\!\text{scattering}}
  = +\rho_0 \!\int \frac{d\kappa_s}{d\Omega\,d\nu}\,
    (\mathbf{p}',\,\mathbf{p})\,I_{\nu'}\,d\Omega'\,d\nu'
    \;-\; \rho_0\,\kappa_s\,I_\nu.
\tag{2.6.21}
\end{equation}
%
By combining equations (2.6.15)--(2.6.21), we obtain our final form of the
equation of radiative transfer
%
\begin{align}
dI_\nu/dl &+ (3\mathbf{n} \cdot \mathbf{a} + \theta
  + 3n^\alpha n^\beta \sigma_{\alpha\beta})\,I_\nu \notag\\
  &= +\rho_0 \!\int \frac{d\kappa_s}{d\Omega\,d\nu}\,
     (\mathbf{p}',\,\mathbf{p})\,I_{\nu'}\,d\Omega'\,d\nu'
     - \rho_0\,\kappa_s\,I_\nu \notag\\
  &\quad + \tfrac{1}{4\pi}\,\rho_0\,\varepsilon_\nu
     - \rho_0\,\kappa_\nu\,I_\nu.
\tag{2.6.22}
\end{align}

%%% -----------------------------------------------------------------
%%% Relationship between emissivity and absorption
%%% -----------------------------------------------------------------

\paragraph*{Relationship between emissivity and absorption}

Into an insulated box place a material medium and a radiation field; and then wait
until complete thermodynamic equilibrium is achieved.\ If $T$ is the equilibrium
temperature, then the radiation field as seen in the LRF will be isotropic with the
standard black-body intensity, $I_\nu = B_\nu$, where
%
\begin{equation}
B_\nu = \frac{2h\nu^3/c^2}{e^{h\nu/kT} - 1}.
\tag{2.6.23}
\end{equation}
%
In equilibrium there must be no net change of $I_\nu$ along any ray: scattering into
the beam must be completely balanced by scattering out of the beam; and
emission must be completely balanced by absorption.\ The requirement of
``detailed balance'' for emission and absorption can be met only if the emissivity
and the absorption coefficient are related by
%
\begin{equation}
\frac{dI_\nu}{dl} = \rho_0\,\kappa_\nu
  \!\left(\frac{\varepsilon_\nu}{4\pi\kappa_\nu}
  - I_\nu\right) = 0;
\tag{2.6.24}
\end{equation}
%
i.e.,
%
\[
\varepsilon_\nu / 4\pi\kappa_\nu = B_\nu.
\]
%
Since $\varepsilon_\nu$ and $\kappa_\nu$ depend only on the thermodynamic state of the matter, and
have nothing to do with the state of the radiation, relation~(2.6.24) must be
satisfied not only inside our insulated box, but also in all other cases where
the matter by itself is in thermodynamic equilibrium but the radiation might
not be.

For matter not in thermodynamic equilibrium one can use a more sophisticated
version of the principle of detailed balance (``Einstein A and B coefficients'') to
derive a more complicated relationship between the emissivity and the absorption
coefficient.\ See, e.g., Chandrasekhar~\cite{Chandrasekhar1960}.

As an application of the equilibrium relationship $\varepsilon_\nu/4\pi\kappa_\nu = B_\nu$, consider free-free
transitions in an ionized gas.\ Whenever the free electrons (which do the
emitting and absorbing) are in thermodynamic equilibrium with each other
(Maxwell--Boltzmann velocity distribution), the absorption coefficient---as
derived from the emissivity (2.1.27)---must be
%
\begin{equation}
\kappa_\nu^{\mathrm{ff}} = (1.50 \times 10^{25}\;\mathrm{cm}^5\;\mathrm{g}^{-1})
  \left(\frac{\rho_0}{\mathcal{A}}\right)
  \left(\frac{(eZ)^2}{\mathcal{A}}\right)
  \bar{G}\;T^{-7/2}\;\nu^{-3}
  \left(\frac{1 - e^{-x}}{x^{-1}}\right),
\tag{2.6.25}
\end{equation}
%
where
\[
x = h\nu/kT.
\]
%
See \S\,2.1 for notation.
